\documentclass[12pt,a4paper]{article}

\usepackage[utf8]{inputenc}
\usepackage[T1]{fontenc}
\usepackage[spanish,es-nodecimaldot]{babel}
\usepackage{amsmath}
\usepackage{amssymb}
\usepackage{amsthm}
\usepackage{geometry}
\usepackage{array}
\usepackage{booktabs}
\usepackage{float}
\usepackage{xcolor}
\usepackage{listings}
\usepackage{graphicx}
\usepackage{enumitem}
\usepackage{tikz}
\usepackage{hyperref}

\begin{document}
\thispagestyle{empty}

\begin{center}
\includegraphics[width=1\textwidth]{encabezado_utvam (1).png}
\vspace{0.2cm}
{\small CÁLCULO IV}
\end{center}

\vspace{2cm}
\begin{center}
{\Large \textbf{Funciones de varias variables}}
\end{center}

\vspace{2.5cm}
\begin{center}
{\Large \textbf{Ejercicios}}
\end{center}

\vspace{3cm}
\begin{center}
{\small \textbf{INTEGRANTES DEL EQUIPO:}}
\end{center}

\vspace{0.1cm}
\begin{center}
{\small 253110061 - López Lemus Glynnes Eileen}
\end{center}

\vspace{0.1cm}
\begin{center}
{\small 253110509 - Rodriguez Martinez Carlos Eduardo}
\end{center}

\vspace{0.1cm}
\begin{center}
{\small 253110366 - Salinas Sanchez Javier}
\end{center}

\vspace{0.1cm}
\begin{center}
{\small 253110017 - Roque Ubaldo Valeria}
\end{center}

\vspace{0.1cm}

\begin{center}
{\small \textbf{GRUPO: TID-401}}
\end{center}

\vspace{0.1cm}
\begin{center}
{\small \textbf{Docente: Rubén Ramírez Gómez}}
\end{center}

\vfill
\begin{flushright}
{\small SEPTIEMBRE, 2026}
\end{flushright}

\newpage

\section{Introducción}

El presente documento contiene el análisis y desarrollo de diferentes
problemas relacionados con funciones de varias variables, superficies
en el espacio tridimensional, vectores, ecuaciones paramétricas y
representación gráfica mediante Wolfram Mathematica.

El trabajo se divide en cuatro problemas principales:

\begin{enumerate}
    \item Optimización de costos en servidores ``Serverless''.
    \item Raycasting manual y determinación del punto de impacto.
    \item Movimiento de un personaje entre dos puntos del espacio.
    \item Reconstrucción matemática de una sonda espacial mediante
    superficies cuadráticas.
\end{enumerate}

En cada problema se muestran las fórmulas utilizadas, las sustituciones,
las operaciones algebraicas y los resultados obtenidos.


\section{Optimización de Costos en Servidores ``Serverless''}

\subsection{Planteamiento del problema}

Una empresa de desarrollo de software utiliza funciones Serverless,
como AWS Lambda, para procesar imágenes.

El costo mensual depende de dos factores:

\begin{itemize}
    \item $n$: número de ejecuciones, expresado en millones.
    \item $t$: tiempo promedio de ejecución por proceso, expresado en segundos.
\end{itemize}

La función proporcionada en el problema es:

\[\boxed{C(n,t)=0.000016nt}\]

donde $C$ representa el costo mensual en dólares.


\subsection{1. Identificación de variables}

Las variables independientes son:

\[\boxed{n \text{ y } t}\]

Esto se debe a que el número de ejecuciones y el tiempo promedio
pueden cambiar y determinar el costo final.

La variable dependiente es:

\[\boxed{C}\]

porque el costo depende directamente de los valores de $n$ y $t$.

Por lo tanto:

\[\boxed{C=C(n,t)}\]

La respuesta puede expresarse de la siguiente manera:

\begin{quote}
Las variables independientes son $n$, que representa el número de
ejecuciones en millones, y $t$, que representa el tiempo promedio
de ejecución en segundos. La variable dependiente es $C$, que
representa el costo mensual en dólares, ya que su valor depende
de $n$ y $t$.
\end{quote}


\subsection{2. Determinación del dominio}

El problema establece las siguientes restricciones.

Para el tiempo de ejecución:

\[0.1\leq t\leq900\]

Para el número de ejecuciones:

\[n\geq0\]

Por lo tanto, el dominio es:

\[\boxed{D=\left\{(n,t)\in\mathbb{R}^2:n\geq0,\quad0.1\leq t\leq900\right\}}\]

También puede escribirse como:

\[\boxed{D=[0,\infty)\times[0.1,900]}\]


\subsection{3. Determinación del rango}

La función es:

\[C(n,t)=0.000016nt\]

Como:

\[n\geq0\]

y:

\[t>0\]

entonces:

\[C\geq0\]

El costo mínimo se obtiene cuando:

\[n=0\]

Sustituyendo:

\[C(0,t)=0.000016(0)(t)\]

\[C=0\]

Por otra parte, el problema no establece un límite superior para $n$.
Por lo tanto, el costo puede crecer indefinidamente.

Así, el rango es:

\[\boxed{R=[0,\infty)}\]


\subsection{4. Representación tabular}

La tabla proporcionada en el ejercicio contiene los siguientes valores:

\begin{table}[H]
\centering
\caption{Costos mensuales}
\begin{tabular}{c|ccc}
\toprule
$n\backslash t$ & $1$ s & $5$ s & $10$ s\\
\midrule
$10$  & $0.16$ & $0.80$ & $1.60$\\
$100$ & $1.60$ & ? & $16.00$\\
$500$ & ? & $40.00$ & ?\\
\bottomrule
\end{tabular}
\end{table}

Existe una inconsistencia en el enunciado.

La fórmula escrita es:

\[C(n,t)=0.000016nt\]

Sin embargo, los valores proporcionados en la tabla corresponden a:

\[\boxed{C(n,t)=0.016nt}\]

Por ejemplo:

\[0.016(10)(1)=0.16\]

\[0.016(10)(5)=0.80\]

\[0.016(10)(10)=1.60\]

Por lo tanto, para completar la tabla de acuerdo con los valores
que aparecen en el ejercicio se utilizará:

\[\boxed{C(n,t)=0.016nt}\]

\subsubsection{Cálculo de $C(100,5)$}

Utilizamos:

\[C=0.016nt\]

Sustituimos:

\[C=0.016(100)(5)\]

Multiplicamos:

\[100(5)=500\]

Entonces:

\[C=0.016(500)\]

Finalmente:

\[\boxed{C=8.00}\]

\subsubsection{Cálculo de $C(500,1)$}

\[C=0.016(500)(1)\]

\[C=0.016(500)\]

\[\boxed{C=8.00}\]


\subsubsection{Cálculo de $C(500,10)$}

\[C=0.016(500)(10)\]

Primero:

\[500(10)=5000\]

Después:

\[C=0.016(5000)\]

Por lo tanto:

\[\boxed{C=80.00}\]


\subsection{Tabla completa}

La tabla completa queda:

\begin{table}[H]
\centering
\caption{Tabla de costos completa}
\begin{tabular}{c|ccc}
\toprule
$n\backslash t$ & $1$ s & $5$ s & $10$ s\\
\midrule
$10$  & $0.16$ & $0.80$ & $1.60$\\
$100$ & $1.60$ & $8.00$ & $16.00$\\
$500$ & $8.00$ & $40.00$ & $80.00$\\
\bottomrule
\end{tabular}
\end{table}

Por lo tanto, los valores faltantes son:

\[\boxed{8.00,\quad 8.00,\quad 80.00}\]

\subsection{5. Análisis gráfico y conceptual}

El problema pregunta qué sucede si el tiempo de ejecución se reduce
a la mitad.

El tiempo nuevo sería:

\[t_{\text{nuevo}}=\frac{t}{2}\]

La función original es:

\[C=kn t\]

donde $k$ representa la constante de proporcionalidad.

Sustituyendo el nuevo tiempo:

\[C_{\text{nuevo}}=kn\left(\frac{t}{2}\right)\]

Por lo tanto:

\[C_{\text{nuevo}}=\frac{k}{2}nt\]

Esto significa que el costo también se reduce a la mitad.

Gráficamente, la superficie presenta una menor pendiente respecto
al eje $n$.

Por lo tanto, la respuesta correcta es:

\[\boxed{\text{B) La superficie se vuelve más plana respecto al eje }n}\]


\section{Raycasting Manual -- El Impacto del Proyectil}

\subsection{Planteamiento}

El proyectil parte del origen:

\[O=(0,0,0)\]

y se dispara en dirección al punto:

\[P=(2,2,2)\]

La pared está determinada por el plano:

\[2x+2y+x=4\]

En este documento se utilizará también la forma:

\[2x+2y+z=4\]

ya que es la forma habitual de un plano tridimensional y produce
el mismo punto de impacto para la trayectoria considerada.


\subsection{1. Vector dirección}

El vector dirección se obtiene mediante:

\[\vec{v}=P-O\]

Sustituyendo:

\[\vec{v}=(2,2,2)-(0,0,0)\]

Restamos componente por componente:

\[\vec{v}=(2-0,2-0,2-0)\]

Por lo tanto:

\[\boxed{\vec{v}=(2,2,2)}\]

\subsection{2. Ecuación paramétrica de la trayectoria}

La ecuación de una recta en el espacio es:

\[\vec{r}(\lambda)=\vec{r}_0+\lambda\vec{v}\]

Como:

\[\vec{r}_0=(0,0,0)\]

y:

\[\vec{v}=(2,2,2)\]

entonces:

\[\vec{r}(\lambda)=(0,0,0)+\lambda(2,2,2)\]

Por lo tanto:

\[\boxed{\vec{r}(\lambda)=(2\lambda,2\lambda,2\lambda)}\]

En forma de ecuaciones:

\[\boxed{
\begin{cases}
x=2\lambda\\
y=2\lambda\\
z=2\lambda
\end{cases}
}
\]


\subsection{3. Cálculo del punto de impacto}

Utilizamos el plano:

\[2x+2y+z=4\]

Sustituimos:

\[x=2\lambda\]

\[y=2\lambda\]

\[z=2\lambda\]

Entonces:

\[2(2\lambda)+2(2\lambda)+(2\lambda)=4\]

Desarrollando:

\[4\lambda+4\lambda+2\lambda=4\]

Sumando:

\[10\lambda=4\]

Despejamos:

\[\lambda=\frac{4}{10}\]

\[\boxed{\lambda=0.4}\]

Ahora sustituimos en las ecuaciones de la trayectoria.

Para $x$:

\[x=2(0.4)=0.8\]

Para $y$:

\[y=2(0.4)=0.8\]

Para $z$:

\[z=2(0.4)=0.8\]

Por lo tanto, el punto de impacto es:

\[\boxed{I=(0.8,0.8,0.8)}\]

También puede expresarse como:

\[\boxed{I=\left(\frac45,\frac45,\frac45\right)}\]

\subsection{4. Comprobación}

Sustituimos el punto obtenido en el plano:

\[2(0.8)+2(0.8)+0.8\]

\[=1.6+1.6+0.8\]

\[=4\]

Por lo tanto:

\[\boxed{2x+2y+z=4}\]

y se confirma que el punto pertenece al plano.


\section{Movimiento de un Personaje}

\subsection{Planteamiento}

El personaje comienza en:

\[A=(2,1,0)\]

y debe desplazarse hasta:

\[B=(5,7,3)\]

\subsection{1. Vector de movimiento}

El vector que va de $A$ hacia $B$ se obtiene mediante:

\[\vec{AB}=B-A\]

Sustituyendo:

\[\vec{AB}=(5,7,3)-(2,1,0)\]

Restamos:

\[\vec{AB}=(5-2,7-1,3-0)\]

\[\boxed{\vec{AB}=(3,6,3)}\]

\subsection{2. Ecuación paramétrica de la trayectoria}

Utilizamos:

\[\vec{r}(t)=A+t(B-A)\]

Sustituyendo:

\[\vec{r}(t)=(2,1,0)+t(3,6,3)\]

Distribuyendo $t$:

\[\boxed{\vec{r}(t)=(2+3t,1+6t,3t)}\]

Por componentes:

\[\boxed{
\begin{cases}
x=2+3t\\
y=1+6t\\
z=3t
\end{cases}
}\]

Para recorrer únicamente el segmento entre $A$ y $B$:

\[\boxed{0\leq t\leq1}\]

\subsection{3. Comprobación de los extremos}

Cuando:

\[t=0\]

obtenemos:

\[x=2+3(0)=2\]

\[y=1+6(0)=1\]

\[z=3(0)=0\]

Por lo tanto:

\[\vec{r}(0)=(2,1,0)=A\]

Cuando:

\[t=1\]

obtenemos:

\[x=2+3=5\]

\[y=1+6=7\]

\[z=3\]

Por lo tanto:

\[\vec{r}(1)=(5,7,3)=B\]

Esto confirma que la trayectoria es correcta.

\subsection{4. Posición después de recorrer el 80\%}

El porcentaje recorrido es:

\[80\%=0.80\]

Por lo tanto:

\[t=0.8\]

Usamos:

\[\vec{r}(t)=(2+3t,1+6t,3t)\]


\subsubsection{Coordenada $x$}

\[x=2+3(0.8)\]

\[x=2+2.4\]

\[\boxed{x=4.4}\]

\subsubsection{Coordenada $y$}

\[y=1+6(0.8)\]

\[y=1+4.8\]

\[\boxed{y=5.8}\]

\subsubsection{Coordenada $z$}

\[z=3(0.8)\]

\[\boxed{z=2.4}\]

Por lo tanto, después de recorrer el 80\% del camino:

\[\boxed{P=(4.4,5.8,2.4)}\]

\section{Operación Reconstrucción: Sonda Espacial Corrupta}

\subsection{Objetivo}

Se tienen las ecuaciones generales de cuatro partes de una sonda
espacial.

El objetivo es transformar cada ecuación a su forma estándar,
identificar la superficie correspondiente y determinar su centro
o vértice.

Las cuatro superficies son:

\begin{enumerate}
    \item Módulo de Energía (Base).
    \item Cápsula de Tripulación (Centro).
    \item Antena de Comunicaciones (Superior).
    \item Núcleo de Procesamiento (Punta).
\end{enumerate}

\subsection{1. Módulo de Energía -- Base}

La ecuación original es:

\[x^2+y^2-z^2+2z-2=0\]

Agrupamos los términos relacionados con $z$:

\[x^2+y^2-(z^2-2z)-2=0\]

Completamos el cuadrado.

Tenemos:

\[z^2-2z\]

Tomamos la mitad de $-2$:

\[\frac{-2}{2}=-1\]

Elevamos al cuadrado:

\[(-1)^2=1\]

Por lo tanto:

\[z^2-2z=(z-1)^2-1\]

Sustituimos:

\[x^2+y^2-\left[(z-1)^2-1\right]-2=0\]

Distribuimos el signo negativo:

\[x^2+y^2-(z-1)^2+1-2=0\]

Simplificamos:

\[x^2+y^2-(z-1)^2-1=0\]

Pasamos $-1$ al segundo miembro:

\[\boxed{x^2+y^2-(z-1)^2=1}\]

La forma estándar también puede escribirse como:

\[\boxed{\frac{x^2}{1}+\frac{y^2}{1}-\frac{(z-1)^2}{1}=1}\]

Como hay dos términos positivos y uno negativo, se trata de:

\[\boxed{\text{Hiperboloide de una hoja}}\]

Su centro es:

\[\boxed{C=(0,0,1)}\]

\subsection{2. Cápsula de Tripulación -- Centro}

La ecuación es:

\[4x^2+4y^2+z^2-8z+12=0\]

Agrupamos:

\[4x^2+4y^2+(z^2-8z)+12=0\]

Completamos el cuadrado en $z$.

Tomamos:

\[\frac{-8}{2}=-4\]

Elevamos al cuadrado:

\[(-4)^2=16\]

Entonces:

\[z^2-8z=(z-4)^2-16\]

Sustituimos:

\[4x^2+4y^2+(z-4)^2-16+12=0\]

Simplificamos:

\[4x^2+4y^2+(z-4)^2-4=0\]

Pasamos $-4$ al segundo miembro:

\[4x^2+4y^2+(z-4)^2=4\]

Dividimos entre $4$:

\[\boxed{x^2+y^2+\frac{(z-4)^2}{4}=1}\]

También:

\[\boxed{\frac{x^2}{1^2}+\frac{y^2}{1^2}+\frac{(z-4)^2}{2^2}=1}\]

Todos los términos son positivos, por lo tanto se trata de:

\[\boxed{\text{Elipsoide}}\]

Su centro es:

\[\boxed{C=(0,0,4)}\]

Sus semiejes son:

\[a=1,\qquad b=1,\qquad c=2\]

El rango de $z$ es:

\[4-2\leq z\leq4+2\]

Por lo tanto:

\[\boxed{2\leq z\leq6}\]

\subsection{3. Antena de Comunicaciones -- Superior}

La ecuación es:

\[x^2+y^2-z+6=0\]

Despejamos $z$:

\[-z=-x^2-y^2-6\]

Multiplicamos por $-1$:

\[\boxed{z=x^2+y^2+6}\]

Esta ecuación corresponde a un:

\[\boxed{\text{Paraboloide circular}}\]

El vértice se obtiene cuando:

\[x=0,\qquad y=0\]

Entonces:

\[z=0^2+0^2+6\]

\[z=6\]

Por lo tanto:

\[\boxed{V=(0,0,6)}\]

La superficie abre hacia:

\[\boxed{+z}\]

y su rango es:

\[\boxed{z\geq6}\]

\subsection{4. Núcleo de Procesamiento -- Punta}

La ecuación es:

\[x^2+y^2+z^2-16z+63.75=0\]

Agrupamos los términos de $z$:

\[x^2+y^2+(z^2-16z)+63.75=0\]

Completamos el cuadrado.

Tomamos:

\[\frac{-16}{2}=-8\]

Elevamos al cuadrado:

\[(-8)^2=64\]

Entonces:

\[z^2-16z=(z-8)^2-64\]

Sustituimos:

\[x^2+y^2+(z-8)^2-64+63.75=0\]

Calculamos:

\[-64+63.75=-0.25\]

Por lo tanto:

\[x^2+y^2+(z-8)^2-0.25=0\]

Pasamos $0.25$ al segundo miembro:

\[\boxed{x^2+y^2+(z-8)^2=0.25}\]

Como:

\[0.25=(0.5)^2\]

podemos escribir:

\[\boxed{x^2+y^2+(z-8)^2=(0.5)^2}\]

Esta es la ecuación de una esfera.

Por lo tanto:

\[\boxed{\text{Esfera}}\]

Centro:

\[\boxed{C=(0,0,8)}\]

Radio:

\[\boxed{r=0.5}\]

\subsection{Resumen de las superficies}

\begin{table}[H]
\centering
\caption{Resumen matemático de la sonda}
\begin{tabular}{p{3.2cm} p{5cm} p{3.5cm} p{3cm}}
\toprule
\textbf{Pieza} &
\textbf{Forma estándar} &
\textbf{Superficie} &
\textbf{Centro/Vértice}\\
\midrule

Base &
$x^2+y^2-(z-1)^2=1$ &
Hiperboloide de una hoja &
$(0,0,1)$\\

Cápsula &
$x^2+y^2+\dfrac{(z-4)^2}{4}=1$ &
Elipsoide &
$(0,0,4)$\\

Antena &
$z=x^2+y^2+6$ &
Paraboloide circular &
$(0,0,6)$\\

Núcleo &
$x^2+y^2+(z-8)^2=0.25$ &
Esfera &
$(0,0,8)$\\

\bottomrule
\end{tabular}
\end{table}

% =================================================
\section{Dominios utilizados para la reconstrucción}

Para representar las diferentes piezas mediante \texttt{ContourPlot3D},
se utilizan los siguientes rangos.

\subsection{Base}

\[-3\leq x\leq3\]

\[-3\leq y\leq3\]

\[0\leq z\leq2.5\]

\subsection{Cápsula}

\[-2\leq x\leq2\]

\[-2\leq y\leq2\]

\[2\leq z\leq6\]

\subsection{Antena}

\[-3\leq x\leq3\]

\[-3\leq y\leq3\]

\[6\leq z\leq8\]

\subsection{Núcleo}

\[-1\leq x\leq1\]

\[-1\leq y\leq1\]

\[7.5\leq z\leq8.5\]

\section{Representación en Wolfram Mathematica}

Para reconstruir gráficamente la sonda se utiliza la función:

\[\texttt{ContourPlot3D}\]

Esta función permite representar superficies tridimensionales
a partir de sus ecuaciones.

\subsection{Definición de las ecuaciones}

\begin{lstlisting}
ec_base = x^2 + y^2 - (z - 1)^2 == 1;

ec_capsula = x^2 + y^2 + (z - 4)^2/4 == 1;

ec_antena = z == x^2 + y^2 + 6;

ec_nucleo = x^2 + y^2 + (z - 8)^2 == 0.25;
\end{lstlisting}


\subsection{Gráfica de la base}

\begin{lstlisting}
base = ContourPlot3D[
   ec_base,
   {x, -3, 3},
   {y, -3, 3},
   {z, 0, 2.5},
   ContourStyle -> Directive[GrayLevel[0.3], Opacity[0.9]],
   Mesh -> None
   ];
\end{lstlisting}

\subsection{Gráfica de la cápsula}

\begin{lstlisting}
capsula = ContourPlot3D[
   ec_capsula,
   {x, -2, 2},
   {y, -2, 2},
   {z, 2, 6},
   ContourStyle -> Directive[Cyan, Opacity[0.45]],
   Mesh -> None
   ];
\end{lstlisting}

\subsection{Gráfica de la antena}

\begin{lstlisting}
antena = ContourPlot3D[
   ec_antena,
   {x, -3, 3},
   {y, -3, 3},
   {z, 6, 8},
   ContourStyle -> Directive[Blue, Opacity[0.75]],
   Mesh -> None
   ];
\end{lstlisting}

\subsection{Gráfica del núcleo}

\begin{lstlisting}
nucleo = ContourPlot3D[
   ec_nucleo,
   {x, -1, 1},
   {y, -1, 1},
   {z, 7.5, 8.5},
   ContourStyle -> Directive[Yellow, Opacity[0.95]],
   Mesh -> None
   ];
\end{lstlisting}

\subsection{Mostrar toda la sonda}

\begin{lstlisting}
Show[
 base,
 capsula,
 antena,
 nucleo,
 PlotRange -> All,
 Boxed -> False,
 Axes -> True,
 AxesLabel -> {"X", "Y", "Z"}
]
\end{lstlisting}


\section{Vistas de la sonda desde diferentes ángulos}

El ejercicio solicita realizar capturas desde al menos cuatro
ángulos diferentes.

Para cambiar el ángulo de visualización se utiliza la opción:

\[
\texttt{ViewPoint}
\]

\subsection{Vista 1}

\begin{lstlisting}
Show[
 base,
 capsula,
 antena,
 nucleo,
 PlotRange -> All,
 Boxed -> False,
 Axes -> True,
 AxesLabel -> {"X", "Y", "Z"},
 ViewPoint -> {2, 2, 2}
]
\end{lstlisting}

\textbf{Captura correspondiente:}

\vspace{3cm}


\subsection{Vista 2}

\begin{lstlisting}
Show[
 base,
 capsula,
 antena,
 nucleo,
 PlotRange -> All,
 Boxed -> False,
 Axes -> True,
 AxesLabel -> {"X", "Y", "Z"},
 ViewPoint -> {-2, 2, 2}
]
\end{lstlisting}

\textbf{Captura correspondiente:}

\vspace{3cm}

\subsection{Vista 3}

\begin{lstlisting}
Show[
 base,
 capsula,
 antena,
 nucleo,
 PlotRange -> All,
 Boxed -> False,
 Axes -> True,
 AxesLabel -> {"X", "Y", "Z"},
 ViewPoint -> {2, -2, 2}
]
\end{lstlisting}

\textbf{Captura correspondiente:}

\vspace{3cm}

\subsection{Vista 4}

\begin{lstlisting}
Show[
 base,
 capsula,
 antena,
 nucleo,
 PlotRange -> All,
 Boxed -> False,
 Axes -> True,
 AxesLabel -> {"X", "Y", "Z"},
 ViewPoint -> {-2, -2, 2}
]
\end{lstlisting}

\textbf{Captura correspondiente:}

\vspace{3cm}

% =================================================
\section{Análisis conceptual de la reconstrucción}

La reconstrucción de la sonda se realizó utilizando las propiedades
de las superficies cuadráticas.

La primera ecuación:

\[x^2+y^2-(z-1)^2=1\]

contiene dos términos positivos y uno negativo, por lo que representa
un hiperboloide de una hoja.

La segunda:

\[x^2+y^2+\frac{(z-4)^2}{4}=\]

contiene tres términos positivos y representa un elipsoide.

La tercera:

\[z=x^2+y^2+6\]

representa un paraboloide circular cuyo vértice se encuentra en
$(0,0,6)$.

Finalmente:

\[x^2+y^2+(z-8)^2=0.25\]

representa una esfera con centro en $(0,0,8)$ y radio $0.5$.

Estas superficies pueden combinarse para representar las diferentes
partes de la sonda espacial.


\section{Resultados finales}

\subsection{Optimización Serverless}

\[\boxed{D=[0,\infty)\times[0.1,900]}\]

\[\boxed{R=[0,\infty)}\]

Valores faltantes de la tabla:

\[\boxed{8.00,\quad8.00,\quad80.00}\]

Si el tiempo se reduce a la mitad, la superficie se vuelve más plana
respecto al eje $n$.

\[\boxed{\text{Respuesta B}}\]

\subsection{Raycasting}

Vector dirección:

\[\boxed{(2,2,2)}\]

Trayectoria:

\[\boxed{\vec{r}(\lambda)=(2\lambda,2\lambda,2\lambda)}\]

Parámetro de impacto:

\[\boxed{\lambda=0.4}\]

Punto de impacto:

\[\boxed{(0.8,0.8,0.8)}\]

\subsection{Movimiento del personaje}

Vector:

\[\boxed{\vec{AB}=(3,6,3)}\]

Trayectoria:

\[\boxed{\vec{r}(t)=(2+3t,1+6t,3t)}\]

Para el 80\%:

\[\boxed{t=0.8}\]

Posición:

\[\boxed{(4.4,5.8,2.4)}\]

\subsection{Sonda espacial}

\begin{align*}
\text{Base:}\quad
&\boxed{x^2+y^2-(z-1)^2=1}\\[0.3cm]
\text{Cápsula:}\quad
&\boxed{x^2+y^2+\frac{(z-4)^2}{4}=1}\\[0.3cm]
\text{Antena:}\quad
&\boxed{z=x^2+y^2+6}\\[0.3cm]
\text{Núcleo:}\quad
&\boxed{x^2+y^2+(z-8)^2=0.25}
\end{align*}

Clasificación:

\begin{itemize}
    \item Base: hiperboloide de una hoja.
    \item Cápsula: elipsoide.
    \item Antena: paraboloide circular.
    \item Núcleo: esfera.
\end{itemize}


\section{Conclusiones}

A través de los ejercicios desarrollados se aplicaron diferentes
conceptos de cálculo multivariable y geometría analítica en tres
dimensiones.

En el problema de optimización de costos se identificaron variables
independientes y dependientes, además de determinar el dominio y
rango de una función de dos variables. También se observó que una
reducción en el tiempo de ejecución produce una reducción proporcional
en el costo.

En el problema de Raycasting se utilizó un vector dirección y una
ecuación paramétrica para determinar el punto en el que una trayectoria
intersecta un plano.

Para el movimiento del personaje se utilizó nuevamente una ecuación
paramétrica de una recta, permitiendo encontrar la posición del
personaje después de recorrer un porcentaje determinado de la
trayectoria.

Finalmente, en la reconstrucción de la sonda espacial se utilizó
el método de completar el cuadrado para transformar ecuaciones
cuadráticas a su forma estándar. Esto permitió identificar un
hiperboloide de una hoja, un elipsoide, un paraboloide circular y una
esfera.

La representación mediante Wolfram Mathematica permitió visualizar
las superficies tridimensionales y combinarlas para reconstruir
matemáticamente la estructura de la sonda.


\section{Observaciones sobre el enunciado}

Se detectaron dos inconsistencias en el material proporcionado.

\subsection{Primera observación: función de costo}

La función indicada en el problema es:

\[C(n,t)=0.000016nt\]

Sin embargo, los valores incluidos en la tabla corresponden a:

\[C(n,t)=0.016nt\]

Por ejemplo:

\[0.016(10)(1)=0.16\]

Por esta razón, para completar la tabla se utilizó $0.016nt$,
manteniendo consistencia con los valores que aparecen en el ejercicio.

\subsection{Segunda observación: ecuación del plano}

El enunciado presenta:

\[2x+2y+x=4\]

Para un plano en tres dimensiones, la forma esperada normalmente
sería:

\[2x+2y+z=4\]

Al utilizar la trayectoria:

\[(x,y,z)=(2\lambda,2\lambda,2\lambda)\]

ambas expresiones conducen al mismo valor:

\[\lambda=0.4\]

y al mismo punto de impacto:

\[\boxed{(0.8,0.8,0.8)}\]

Por lo tanto, el resultado del punto de impacto se mantiene.


\end{document}
